\documentclass[11pt,a4paper]{article}
\usepackage[utf8]{inputenc}
\usepackage[T1]{fontenc}
\usepackage[spanish]{babel}
\usepackage{hyperref}
\usepackage{geometry}
\geometry{margin=2.5cm}
\usepackage{parskip}

\title{103 sin archivo no hay resultado}
\author{Alejandro Valenzuela Rivera\\\small Universidad Autónoma de Nuevo León}
\date{\today}
\begin{document}
\maketitle

\begin{quote}
\textbf{Sesgo declarado.} Este trabajo se sostiene desde el sesgo declarado de su autor, y no finge neutralidad. Lo que afirma se puede tumbar, y ese es el punto: el falsificador está escrito, no escondido.
\end{quote}

\section*{Resumen}
Un resultado solo existe si hay un archivo con lo que se corrio y su salida. Si no hay archivo, la seccion dice sin datos, y sin datos es un resultado: el unico que se puede escribir sin medir nada.

\section{Por qué se sostiene}
Y es la respuesta a no alucinar resultados: la alucinacion no es un error, es reflejo; el error es no decir de donde viene. Un resultado que no se puede citar con su archivo no tiene de donde venir, y por eso se lee como inventado aunque no lo sea.

\section{De dónde viene}
La autoridad, 2026-10-25.

\section{El caso}
Ninguna: es la regla de la seccion.

\section{Método}
Un resultado con su archivo, y con la corrida que lo produjo.

\section{Límites, y qué lo tumba}
Un resultado sin archivo: una historia que se puede contar de cualquier forma.

\section{Lo que queda abierto}
\textquestiondown{}Que se escribe cuando no se midio nada?

\section*{Cómo se leyó esta tesis}
Esto se escribió en el repositorio man-made-empathy, del que este candado es una pieza: candados/04 el lenguaje/103 sin archivo no hay resultado.md. La máquina que lo compuso lo hizo desde este mismo archivo, y no agregó nada que no estuviera escrito. Si algo de esto es falso, se dice en el mismo sitio, y se corrige con fecha.

\end{document}
